\documentclass[a4paper, serif, 10pt]{moderncv}

%% moderncv
% color and style
\moderncvcolor{black}
\moderncvstyle{banking}

%% page layout/margins
\usepackage[top=2.5cm, bottom=2.5cm, left=1.5cm, right=1.5cm]{geometry}

\nopagenumbers{}

%% Babel config for Italian language
% ref:
% - https://latex3.github.io/babel/guides/locale-italian.html
\usepackage[italian]{babel}

%% fontspec
\usepackage{fontspec}

\IfFontExistsTF{Linux Libertine}{
  \setmainfont[Ligatures={TeX, Common}, Numbers=OldStyle]{Linux Libertine}
}{}
\IfFontExistsTF{Linux Libertine O}{
  \setmainfont[Ligatures={TeX, Common}, Numbers=OldStyle]{Linux Libertine O}
}{}

%% CTeX
% CJK support, used to show CJK characters in the resume
%
% - heading=false: disable ctex heading
% - punct=kaiming: use kaiming punctuations styles for CJK
% - scheme=plain: use plain scheme, do not override \normalsize font size
% - space=auto: space settings for CJK characters
%
% CJK fonts are resolved per locale language, e.g. Noto Serif/Sans CJK JP
% for Japanese resumes, with fallback to the next region if not installed
%
% ref:
% - http://ctan.mirrorcatalogs.com/language/chinese/ctex/ctex.pdf
\usepackage[UTF8, heading=false, punct=kaiming, scheme=plain, space=auto]{ctex}

\IfFontExistsTF{Noto Serif CJK SC}{
  \setCJKmainfont{Noto Serif CJK SC}
}{}
\IfFontExistsTF{Noto Sans CJK SC}{
  \setCJKsansfont{Noto Sans CJK SC}
}{}

%% line spacing
\usepackage{setspace}
\setstretch{1.125}

%% URL styling - use normal text instead of monospace
\urlstyle{same}

%% Patch \httplink and \httpslink to handle full URLs
% The original moderncv macros always prepend http:// or https:// to URLs.
% This causes issues when the URL already has a protocol (e.g., https://example.com)
% resulting in malformed URLs like https://https://example.com.
% This patch checks if the URL already contains "://" and uses it directly if so.
% Note: We use \str_set:Nx (x-expansion) to expand macros like \@homepage first.
\ExplSyntaxOn
\str_new:N \g_yamlresume_url_str

\RenewDocumentCommand{\httpslink}{O{}m}{
  \str_set:Nx \g_yamlresume_url_str {#2}
  \str_if_in:NnTF \g_yamlresume_url_str {://}
    { \url{#2} }
    { \url{https://#2} }
}

\RenewDocumentCommand{\httplink}{O{}m}{
  \str_set:Nx \g_yamlresume_url_str {#2}
  \str_if_in:NnTF \g_yamlresume_url_str {://}
    { \url{#2} }
    { \url{http://#2} }
}
\ExplSyntaxOff

\name{Giulia Ferrari}{}
\title{Data Scientist | Machine Learning e Analytics}
\phone[mobile]{+39 347 123 4567}
\email{giulia.ferrari@example.it}
\homepage{https://www.giuliaferrari-data.it}

\address{Via Giovanni Battista Pirelli 18, Milano, Lombardia, Italia, 20124}{}{}

\extrainfo{{\small \faLinkedin}\ \href{https://www.linkedin.com/in/giulia-ferrari-data}{@giulia-ferrari-data} {} {} {} • {} {} {} 
{\small \faGithub}\ \href{https://github.com/giuliaferrari}{@giuliaferrari} {} {} {} • {} {} {} 
{\small \faStackOverflow}\ \href{https://stackoverflow.com/users/1234567/giulia-ferrari}{@giulia-ferrari}}

\begin{document}

\maketitle

\section{Informazioni personali}

\cvline{}{Data Scientist con oltre 7 anni di esperienza nella progettazione di modelli predittivi e soluzioni analitiche per i settori energia, bancario e pagamenti digitali.
%
\begin{itemize}
\item Esperienza end-to-end: dalla raccolta dei requisiti alla messa in produzione e al monitoraggio di modelli.
\item Solida base statistica e di apprendimento automatico, con attenzione alla validazione, all'interpretabilità e all'impatto di business.
\item Capacità di comunicare risultati complessi a stakeholder tecnici e non tecnici, in italiano e in inglese.
\item Approccio collaborativo, mentoria di junior e forte orientamento alla qualità del dato.
\end{itemize}}
\section{Formazione}

\cventry{set 2015–lug 2017}
        {Laurea magistrale, Data Science, Punteggio: 110/110 e lode}
        {Università degli Studi di Milano-Bicocca}
        {\url{https://www.unimib.it}}
        {}
        {\begin{itemize}
\item Formazione avanzata in statistica, machine learning e gestione di grandi volumi di dati.
\item Sviluppo di progetti applicativi con Python, R e SQL su casi reali di previsione e segmentazione.
\item Particolare attenzione alla validazione dei modelli, alla comunicazione dei risultati e all'impatto organizzativo.
\end{itemize}
\textbf{Corsi}: Statistica Computazionale, Apprendimento Automatico, Basi di Dati e Data Warehousing, Analisi delle Serie Temporali, Data Mining, Visualizzazione dei Dati, Metodi Bayesiani, Etica e Governance dei Dati}

\cventry{set 2012–lug 2015}
        {Laurea triennale, Statistica e Informatica, Punteggio: 108/110}
        {Università degli Studi di Milano}
        {\url{https://www.unimi.it}}
        {}
        {\begin{itemize}
\item Solide basi quantitativa e informatica per l'analisi dei dati.
\item Esperienza pratica con dataset reali, pulizia dei dati e reportistica statistica.
\end{itemize}
\textbf{Corsi}: Statistica Descrittiva, Probabilità e Inferenza, Programmazione in Python, Basi di Dati Relazionali, Analisi Esplorativa dei Dati, Metodi di Campionamento}
\section{Esperienza professionale}

\cventry{mar 2022–Presente}
        {Data Scientist Senior}
        {Nexi Digital}
        {\url{https://www.nexi.it}}
        {}
        {\begin{itemize}
\item Guido un team di 4 data scientist nello sviluppo di modelli predittivi per il rilevamento delle frodi e la personalizzazione delle offerte.
\item Progettato e implementato pipeline MLOps su AWS per l'addestramento, il versionamento e il monitoraggio continuo dei modelli.
\item Ridotto del 18\% i falsi positivi nel sistema antifrode migliorando le regole di validazione e l'interpretabilità dei modelli.
\item Collaborato con product manager, ingegneri e team compliance per tradurre requisiti di business in soluzioni dati misurabili.
\item Mentorato due junior data scientist e promosso pratiche di code review e documentazione tecnica.
\end{itemize}
\textbf{Parole chiave}: Modelli predittivi, MLOps, Customer Analytics, Leadership tecnica}

\cventry{gen 2019–feb 2022}
        {Data Scientist}
        {Enel X}
        {\url{https://www.enelx.com/it/it}}
        {}
        {\begin{itemize}
\item Sviluppato modelli di previsione della domanda energetica a livello giornaliero e orario, con riduzione dell'errore MAPE dal 12\% all'8\%.
\item Realizzato segmentazioni clienti tramite clustering e analisi di sopravvivenza per campagne di retention mirate.
\item Automatizzato la reportistica mensile con Airflow e Power BI, riducendo il tempo di preparazione da 3 giorni a 4 ore.
\item Condotto test A/B e analisi causali per valutare l'impatto delle iniziative commerciali.
\end{itemize}
\textbf{Parole chiave}: Previsione della domanda, Customer Analytics, Test A/B, Airflow}

\cventry{set 2017–gen 2019}
        {Analista Dati}
        {UniCredit}
        {\url{https://www.unicredit.it}}
        {}
        {\begin{itemize}
\item Analizzato portafogli creditizi e costruito dashboard per il monitoraggio degli indicatori di rischio.
\item Sviluppato modelli statistici per la previsione di insolvenza e la valutazione del merito creditizio.
\item Pulito e integrato dati provenienti da fonti eterogenee, migliorando la qualità e la coerenza delle analisi.
\item Presentato risultati periodici ai team risk e business con report chiari e visualizzazioni efficaci.
\end{itemize}
\textbf{Parole chiave}: Rischio di credito, SQL, Dashboard, Statistica}
\section{Lingue}

\cvline{Italiano}{Madrelingua o bilingue \hfill \textbf{Parole chiave}: Madrelingua}
\cvline{Inglese}{Competenza professionale completa \hfill \textbf{Parole chiave}: IELTS 8.0, C1}
\cvline{Francese}{Competenza professionale limitata \hfill \textbf{Parole chiave}: DELF B1}
\section{Competenze}

\cvline{Apprendimento Automatico}{Esperto \hfill \textbf{Parole chiave}: Python, scikit-learn, TensorFlow, PyTorch, XGBoost}
\cvline{Ingegneria dei Dati}{Avanzato \hfill \textbf{Parole chiave}: SQL, Spark, Airflow, dbt, Databricks}
\cvline{Statistica Applicata}{Esperto \hfill \textbf{Parole chiave}: Test A/B, Inferenza causale, Serie temporali, Analisi di sopravvivenza}
\cvline{Visualizzazione dei Dati}{Avanzato \hfill \textbf{Parole chiave}: Power BI, Tableau, Looker, Plotly, Streamlit}
\cvline{Cloud e MLOps}{Avanzato \hfill \textbf{Parole chiave}: AWS, Azure, Docker, Kubernetes, MLflow}
\cvline{Comunicazione dei Risultati}{Avanzato \hfill \textbf{Parole chiave}: Narrazione dei dati, Presentazioni executive, Documentazione tecnica, Facilitazione workshop}
\section{Premi}

\cventry{giu 2023}
        {Associazione Italiana per l'Intelligenza Artificiale}
        {Premio Innovazione Dati 2023}
        {}
        {}
        {Riconoscimento per il progetto "Previsione della Domanda Energetica", selezionato tra le migliori applicazioni di intelligenza artificiale in ambito utilities.}

\cventry{ott 2016}
        {Università degli Studi di Milano-Bicocca}
        {Borsa di Studio per Merito}
        {}
        {}
        {Borsa di studio assegnata per eccellenza accademica durante il corso di laurea magistrale in Data Science.}
\section{Certificazioni}

\cventry{mag 2021}
        {Google Cloud}
        {Google Professional Data Engineer}
        {\url{https://cloud.google.com/certification/data-engineer}}
        {}
        {}

\cventry{nov 2020}
        {Amazon Web Services}
        {AWS Certified Machine Learning – Specialty}
        {\url{https://aws.amazon.com/certification/}}
        {}
        {}

\cventry{set 2022}
        {Microsoft}
        {Microsoft Certified: Azure Data Scientist Associate}
        {\url{https://learn.microsoft.com/en-us/certifications/azure-data-scientist/}}
        {}
        {}
\section{Pubblicazioni}

\cventry{dic 2022}
        {Modelli predittivi per la manutenzione di reti elettriche intelligenti}
        {Rivista Italiana di Statistica Applicata}
        {\url{https://doi.org/10.1234/risa.2022.12.004}}
        {}
        {\begin{itemize}
\item Confronto tra modelli di serie temporali e apprendimento supervisionato per la manutenzione predittiva.
\item Analisi dell'impatto delle variabili meteo e di carico sulla probabilità di guasto.
\end{itemize}}

\cventry{set 2021}
        {Segmentazione dei clienti tramite apprendimento non supervisionato}
        {Atti della Conferenza Italiana di Data Science}
        {\url{https://doi.org/10.5678/cids.2021.09.012}}
        {}
        {\begin{itemize}
\item Applicazione di clustering e riduzione della dimensionalità a dati di consumo energetico.
\item Valutazione della stabilità dei segmenti e delle implicazioni per campagne di marketing.
\end{itemize}}
\section{Progetti}

\cventry{mar 2021–dic 2021}
        {Sistema di previsione della domanda energetica per un'azienda utilities, con modelli a livello giornaliero e orario.}
        {Previsione della Domanda Energetica}
        {\url{https://github.com/giuliaferrari/energy-demand-forecast}}
        {}
        {\begin{itemize}
\item Costruito pipeline di feature engineering con dati storici, meteo e calendario.
\item Confrontato Prophet, XGBoost e LSTM; selezionato il modello con il miglior errore in validazione temporale.
\item Pubblicato dashboard interattiva per il monitoraggio delle previsioni e degli scostamenti.
\end{itemize}
\textbf{Parole chiave}: Serie temporali, Prophet, XGBoost, Airflow, Dashboard}

\cventry{giu 2020–nov 2020}
        {Analisi di recensioni e sondaggi per identificare i fattori chiave della soddisfazione clienti.}
        {Analisi della Soddisfazione Clienti}
        {\url{https://github.com/giuliaferrari/customer-satisfaction}}
        {}
        {\begin{itemize}
\item Applicato NLP per estrarre topic e sentiment da recensioni in italiano.
\item Integrato i risultati con dati transazionali per stimare l'impatto sulla soddisfazione.
\item Presentato insight actionable al team customer experience.
\end{itemize}
\textbf{Parole chiave}: NLP, Topic modeling, Sentiment analysis, Power BI}
\section{Interessi}

\cvline{Sport}{Corsa, Sci, Yoga}
\cvline{Musica}{Pianoforte, Jazz}
\cvline{Divulgazione}{Mentoraggio, Open data, Scrittura tecnica}
\section{Volontariato}

\cventry{gen 2020–giu 2024}
        {Mentor Data Science}
        {Dati per il Bene Comune}
        {\url{https://www.datiperilbenecomune.it}}
        {}
        {\begin{itemize}
\item Supportato associazioni no profit nella definizione di indicatori e dashboard per misurare l'impatto sociale.
\item Tenuto laboratori introduttivi su Python, analisi dati e visualizzazione per volontari e studenti.
\item Contribuito alla pubblicazione di dataset aperti e linee guida per la qualità del dato.
\end{itemize}}
    
\end{document}